\songtitle{Bogotá Bajo Tierra}

\tag*{2026}\tag{tropipop}\tag{thomas-entourage}\tag{infrastructure}\tag{spanish-lyrics}
\begin{notes}
Lyrics by SUNO based on a Gemini Notebook summary of a document co-written by Carlos Thompson and Claude, from Carlos Thompson's original works adapted to \emph{Thomas’ Entourage} narrative universe.
\end{notes}
\begin{style}
Colombian tropipop jingle with upbeat guacharaca groove and syncopated percussion; verse rides light acoustic guitar and marimba-like synth pulses, pre-chorus strips to handclaps and bass lift, chorus opens with bright brass hits and crowd-ready hooks, Warm lead vocal with tight doubles, playful call-and-response backing lines, short delay throws on key words, risers into each chorus, and a glossy, sunny mix that feels urban and hopeful
\end{style}
\begin{lyrics}
\songpart{Intro}
Bogotá, Bogotá

\songpart{Verse 1}
InfraVia lo trazó\\
en el año dos mil veintiséis\\
un plan que mira lejos\\
para hacer respirar la ciudad

Bajo el asfalto\\
van a correr los largos viajes\\
y arriba, poco a poco\\
la calle vuelve a caminar

\songpart{Pre-Chorus}
Si el peso baja al túnel\\
la superficie se despierta\\
más paso, más bici\\
más barrio, más puerta

\songpart{Chorus}
Bogotá bajo tierra
(¡bajo tierra!)\\
Bogotá más abierta
(más abierta)\\
túneles para el pulso largo\\
bulevares verdes arriba\\
Bogotá bajo tierra\\
y la vida va más viva

\songpart{Verse 2}
Fases de veinte años\\
con futuro en cada plano\\
unir el metro que viene\\
con los rieles del pasado

Corredores que eran largos\\
de espera, humo y tensión\\
se vuelven senda y parque\\
para cruzar con corazón

\songpart{Pre-Chorus}
Tráfico pesado lejos\\
y la superficie sonríe\\
entre árboles y andenes\\
la ciudad se redefine

\songpart{Chorus}
Bogotá bajo tierra
(¡bajo tierra!)\\
Bogotá más abierta
(más abierta)\\
túneles para el pulso largo\\
bulevares verdes arriba\\
Bogotá bajo tierra\\
y la vida va más viva

\songpart{Bridge}
Una ciudad generacional\\
que no se queda quieta\\
ingenio en la estructura\\
y habitabilidad en la meta

Metro, bici y peatón\\
en la misma melodía\\
Bogotá se transforma\\
con eficiencia y armonía

\songpart{Chorus}
Bogotá bajo tierra
(¡bajo tierra!)\\
Bogotá más abierta
(más abierta)\\
túneles para el pulso largo\\
bulevares verdes arriba\\
Bogotá bajo tierra\\
y la vida va más viva

\songpart{Outro}
Bogotá, Bogotá\\
más habitable cada día
\end{lyrics}

